\section{Equivariant cones and an lc projective slice}\label{sec:cones}

Throughout the local argument the ground field is $\C$.  If $\theta$ is
a nonzero rational volume form, we write $A_\theta$ for the log discrepancy
of the sub-pair whose boundary on a smooth model is
$-\operatorname{div}(\theta)$.  For an $s$-fold pluri-volume form the
boundary is $-s^{-1}\operatorname{div}(\theta)$.  These conventions are
homogeneous in valuations and are compatible with separable finite
extensions: pulling back the form gives the same discrepancy on an
extended valuation, with its actual normalization.  This is the
characteristic-zero ramification formula.  An \emph{lc form} has
nonnegative discrepancy at every divisorial valuation under consideration;
on a projective model this tests the entire function field.  A generator
with zero Weil divisor on a normal variety recovers the usual discrepancy
at centres on that variety.  On real quasi-monomial valuations we use the
log-linear extension on log-smooth models.

Consider a sequence of pointed normal Gorenstein klt germs $(V,o)$ of
fixed dimension $n\geq2$, with a cyclic group fixing $o$ and a
semi-invariant generator $\theta_V$ of zero divisor.  We shrink affine
neighbourhoods equivariantly when necessary.  Constants denoted by $C$
may change, and a uniform assertion about the sequence is allowed to hold
after passage to a subsequence.

\subsection{Rational regrading on the whole local algebra}

The stable degeneration theorem
\cite[Theorem~1.2(1)--(2)]{StableDegeneration} and uniqueness of the
normalized-volume minimizer \cite[Theorem~1.1]{MinimizerUniqueness}
give a quasi-monomial valuation $v_*$, normalized by $A(v_*)=1$, with
finitely generated associated graded algebra $R$. The affine cone
$C_0=\Spec R$ is normal and klt, and its induced Reeb valuation minimizes
normalized volume. The valuation $v_*$ is invariant under the cyclic
group \cite[Corollary~1.2]{MinimizerUniqueness}. Only this first cone
will be used.

We first explain why a nearby rational monomial valuation has this same
graded algebra with a different grading. Related rational regrading is
proved in \cite[Lemma~2.10]{LiXuHigherRank}; here we keep track of the
whole local algebra and the cyclic action explicitly. Present $v_*$ at the generic
point of an SNC stratum on a smooth model.  After a toroidal subdivision
respecting the smallest rational subspace containing its weight vector,
we may work on a face with positive rationally independent weights.
Write $\rho$ for its number of monomial coordinates.  Discrepancy is
linear on this face.  Choose regular semi-invariant lifts
$h_1,\ldots,h_r$ of homogeneous algebra generators of $R$, using character
projections for the finite group.  Enlarge this list by regular
eigenfunctions generating the maximal ideal and, when needed, by centred
affine coordinates.

In the completed monomial chart, with coefficients in a residue
coefficient field, each $h_i$ has one least exponent for the original
weights.  That exponent remains uniquely least for every sufficiently
nearby weight vector.  To justify a common neighbourhood, exponents far
out in $\N^\rho$ have uniformly large weight when all coordinate weights
stay in a compact subset of the positive cone; only finitely many
exponents can compete with the chosen one.  There are finitely many
generators.

Fix one neighbourhood $U$ of positive weight vectors on which every
$h_i$ retains its unique least exponent $\alpha_i$ and its coefficient
$c_i\ne0$. This same neighbourhood works for products of every degree.
Indeed, in $h^\beta=\prod_i h_i^{\beta_i}$, every exponent other than
$E_\beta=\sum_i\beta_i\alpha_i$ differs from $E_\beta$ by a sum of
generator exponent gaps. Each nonzero gap pairs positively with every
weight in $U$, so $E_\beta$ remains uniquely least, with coefficient
$\prod_i c_i^{\beta_i}$. No bound on $\beta$ is needed.

For any regular local function $g$, successive subtraction of its
initial form expresses it, to arbitrarily high $v_*$-order, as a sum of
polynomials $P_j(h)$ homogeneous for the original weights. The process
goes to infinite order because the finitely generated monoid of positive
weights has only finitely many elements in a bounded interval. In one
such polynomial all the exponents $E_\beta$ coincide: their original
scalar weights coincide and the original coordinate weights are
rationally independent. Its coefficient sum at this exponent is nonzero,
because $P_j$ was chosen to represent the nonzero initial of the current
remainder. The product calculation therefore shows that $P_j(h)$ retains
that exponent and coefficient for every weight in the same $U$.
Distinct subtraction groups have distinct exponent vectors, so even if
their new scalar weights tie, their initial monomials cannot cancel.
Finally, on a compact positive subneighbourhood of $U$, coordinatewise
weight comparison gives $v_{\rm new}(r)\ge c v_*(r)$ for every regular
subtraction error $r$, with one $c>0$. The errors thus tend to infinite
order for every such new valuation. Consequently the new associated
graded ring is exactly the subring generated by the same initial
monomials in the polynomial ring over the residue coefficient field.

This conclusion is equivariant on the entire local algebra.  Indeed,
relations and leading cancellations split by their exponent vectors,
and their equivariance is inherited from the $v_*$-filtration.  For
completeness, the error comparison also holds after every group
translation.  A finite subtraction error $r$ is regular at $(V,o)$, as
is $\gamma r$ because $\gamma$ fixes $o$.  Both pull back to regular
functions on the fixed monomial chart, even if $\gamma$ does not preserve
that chart.  Their series therefore have nonnegative boundary
exponents.  For a fixed positive comparison constant $c$,
\[
 v_{\mathrm{new}}(\gamma r)
 \geq c\,v_*(\gamma r)=c\,v_*(r)\longrightarrow\infty.
\]
Only this comparison for regular errors is needed.

Choose sufficiently close positive rational weights and rescale them to
a primitive integral valuation $v_{\mathrm{int}}$.  Since $V$ is
Gorenstein and klt,
\begin{equation}\label{cone:normalization}
 a=A_{\theta_V}(v_{\mathrm{int}})\in\Z_{>0},
 \qquad v=a^{-1}v_{\mathrm{int}},\qquad A_{\theta_V}(v)=1.
\end{equation}
Integral degree $I$ in $R$ thus means normalized degree $I/a$.  We choose
the weights close enough that every positive monomial weight, after
normalization, lies between one half and twice its $v_*$-weight.

The associated extended Rees family is of finite type.  More explicitly,
on an affine neighbourhood containing the chosen generators, the ideal
of functions of $v_{\mathrm{int}}$-order at least $j>0$ is generated by
monomials in the $h_i$ having weight at least $j$.  In the local ring,
successive initial subtraction gives this statement modulo errors of
arbitrarily high valuation.  Izumi's inequality for the divisorial
valuation at the klt germ \cite[Theorem~3.1]{LiMinimizing}, with a
constant depending on this fixed germ and valuation, makes these errors
arbitrarily deep in the maximal ideal; Krull intersection then gives the
ideal equality.  Away from $o$ both ideals are the unit ideal, because
the chosen generators cut out $o$.  The family is therefore generated
by the original algebra, $t$, and
$t^{-v_{\mathrm{int}}(h_i)}h_i$.  It is flat over the $t$-line, has a
section through the pointed germs, and has central fibre $C_0$.
It is normal: away from the central fibre this follows from normality of
$V$, while the normal Cohen--Macaulay Cartier fibre gives the depth and
codimension-one conditions along that fibre.

Order along the central fibre restricts to $v_{\mathrm{int}}$ and is
its Gauss extension with $w(t)=1$.  Indeed, when terms involving distinct
powers of $t$ tie, their coefficient initials belong to distinct degrees
of $R$ and cannot cancel.  The monomial discrepancy formula on the
product with the logarithmic $t$-line consequently shows that
\begin{equation}\label{cone:rees-form}
 t^{-a}\theta_V\wedge\frac{dt}{t}
\end{equation}
has exactly a simple pole along the central fibre and no other divisor.
Thus the total space has Cartier canonical class, and residue along the
smooth locus of the normal central fibre gives a generator $\theta$ on
$C_0$ with zero Weil divisor and integral weight $a$.

This generator is also a character for the original grading torus.
The reason is that units of a positive graded affine domain with
degree-zero part $\C$ are constants.  The generators produced by
different rational regradings therefore differ by scalars.  Their
normalized canonical weights are all one, so continuity of pairing
with the fixed canonical character gives canonical weight one for the
original Reeb valuation.  This weight equals its discrepancy, as in the Reeb-character formula
of \cite[Lemma~2.18 and Definition~2.20]{LiXuHigherRank}. In the present
volume-form convention this can also be seen directly: choose
homogeneous rational monomials forming a basis of the character lattice,
view them as transcendental torus coordinates over the degree-zero
field, and divide the form by its character monomial.  What remains is
a base volume form wedged with logarithmic fibre differentials, to
which the weighted Gauss formula applies.  The same description proves
continuity of discrepancy for these Reeb weights.

\subsection{Alpha estimates and comparison of valuations}

\begin{proposition}\label{prop:cone-estimates}
For the preceding choices, every nonzero homogeneous $h_I\in R_I$ of
positive degree satisfies
\begin{equation}\label{cone:alpha}
 \lct_o(C_0;h_I)\geq\frac{1}{2n(I/a)}.
\end{equation}
Every nonzero regular local function $g$ on $V$ and every divisorial or
real quasi-monomial valuation $w$ centred at $o$ satisfy
\begin{equation}\label{cone:alpha-germ}
 w(g)\leq 2n A_{\theta_V}(w)\,v(g).
\end{equation}
Moreover, for every valuation $w$ centred at $o$,
\begin{equation}\label{cone:comparison}
 w(g)\geq s(w)v(g)\quad(g\in\cO_{V,o}),
 \qquad s(w)=\min_i\frac{w(h_i)}{v(h_i)}.
\end{equation}
In particular each positive normalized homogeneous degree is at least
$1/(2n)$.
\end{proposition}

\begin{proof}
First take an eigenfunction $f$ for the original grading torus, of
positive weight.  An equivariant log resolution gives an invariant
divisorial valuation computing its threshold at the vertex.  Normalize
it to $A(w)=1$ and set $c=w(f)>0$.  It is nonnegative on $R$.
For a nearby rational normalized Reeb grading $\zeta$ and a positive
rational number $s$, define
\[
 u_s(h_\chi)=\zeta(h_\chi)+s w(h_\chi)
\]
on torus eigenvectors and take the minimum over characters on sums.
This is a valuation with
\begin{equation}\label{cone:twisted-discrepancy}
 A(u_s)=1+s.
\end{equation}
Here is a direct verification of the discrepancy statement.  On the
product with a logarithmic line of coordinate $z$, take the Gauss
extension of $sw$ with $w'(z)=\gamma>0$ rational.  Choose an integral
one-parameter subgroup $j$ and $\gamma$ so that
$\zeta(\chi)=\gamma j(\chi)$, and apply the field automorphism
$h_\chi\mapsto h_\chi z^{j(\chi)}$, fixing $z$.  Its restriction has
the displayed values $u_s$.  It is still a Gauss extension: after
clearing denominators, coefficients that tie at a power of $z$ do not
cancel across distinct characters, because the $w$-filtration is torus
invariant.  The transformation of $\theta\wedge d\log z$ adds the
canonical character weight $1$ to discrepancy.  The original Gauss
extension has discrepancy $s$.  Both it and the transformed restriction
are divisorial up to rational scaling: the value group is discrete and
nonzero, and restriction from the product leaves residue transcendence
degree at least $n-1$.  The logarithmic Gauss extension has the same
discrepancy as its restriction.  These observations prove
\eqref{cone:twisted-discrepancy}.

Put $d_\zeta=\zeta(f)$.  Multiplication by $f$ shifts the $u_s$-filtration
exactly, while $u_s\geq\zeta$ on $R$.  Counting modulo consecutive
powers of $f$ therefore gives
\begin{equation}\label{cone:twisted-volume}
 \vol(u_s)\leq\vol(\zeta)
                \frac{d_\zeta}{d_\zeta+sc}.
\end{equation}
Indeed, at cutoff $p$ the length is at most the sum, over $k\geq0$, of
the dimensions in $R/(f)$ up to $\zeta$-weight
$p-k(d_\zeta+sc)$.  The cumulative Hilbert growth of $R/(f)$ has leading
coefficient $d_\zeta\vol(\zeta)/(n-1)!$, by the nonzerodivisor identity
for the rational graded Hilbert series, after integral rescaling.
Summing this polynomial growth proves \eqref{cone:twisted-volume};
local lengths at the vertex equal these affine counts.

Use normalized-volume minimization on $C_0$ and let $\zeta$ approach its
minimizing Reeb valuation.  Multiplicative comparison of weights gives
continuity of volume.  Writing $d_*=v_*(f)$, we obtain
\[
 d_*+sc\leq(1+s)^n d_*.
\]
Letting $s$ decrease to zero yields $c\leq nd_*$.  A homogeneous $h_I$
for the rational regrading may have several components for the original
torus.  A one-parameter torus degeneration selects a nonzero extreme
component.  Semicontinuity of the threshold in the trivial family,
equivalently inversion of adjunction, bounds the threshold of the
general translate below by that of this component.  Its original weight
is at most $2I/a$.  This proves \eqref{cone:alpha}.

For a nonunit $g$, the function
$t^{-v_{\mathrm{int}}(g)}g$ is regular near the special point of the
Rees family and restricts to its graded initial.  Inversion of
adjunction for a normal Cartier divisor in a Gorenstein variety
\cite{KawakitaInversion} transfers the threshold bound to the general
pointed germ.  Its different is restriction here, since the total space
is smooth along the smooth locus of the central fibre.  More explicitly,
the pair obtained by adding the central divisor is lc near the special
point; away from $t=0$ the family of divisors is a product, and the
pointed section gives the desired bound at $o$.  This is exactly
\eqref{cone:alpha-germ} for divisorial valuations.  Approximation on a
log-smooth model gives it for real quasi-monomial valuations.  Units
have value zero and cause no exception.

The monomial generators of valuation ideals proved above give
\eqref{cone:comparison}: each generating monomial of $v$-weight at least
$v(g)$ has $w$-weight at least $s(w)v(g)$.  Finally a nonunit at the
vertex has threshold at most one, so \eqref{cone:alpha} implies the
lower bound $I/a\geq1/(2n)$.
\end{proof}

\subsection{Complements and a root cover}

There is a homogeneous nonzero function
\begin{equation}\label{cone:complement-section}
 f_*\in R_{aN},\qquad
 \left(C_0,\frac1N\operatorname{div}(f_*)\right)\text{ is lc},
\end{equation}
where the positive integer $N$ is bounded in terms of $n$.
No finite-group equivariance of $f_*$ is needed.
To prove this, let $E=\Proj R$ and $H=\cO_E(1)$ as an ample rational
class.  The quotient carries the standard boundary $B$, with coefficient
$1-1/e_P$ at a prime of generic isotropy order $e_P$.  The grading is
effective because $v_{\mathrm{int}}$ is primitive.  On a homogeneous
principal open, invert a homogeneous section $b$ and take the slice
$b=1$.  Its product with the radial torus maps finite \'etale to that
cone open, and its cyclic quotient gives the corresponding Proj open
with the stated ramification boundary.  The slice is klt, so $(E,B)$
is klt.

Choose a homogeneous rational element $s$ of weight one.  Then
$k(C_0)=k(E)(s)$, and write
\[
 \theta=s^a\omega\wedge d\log s.
\]
The slice description identifies the divisor of the logarithmic product
form over every prime of $E$ with the ramification multiple of
$K_E(\omega)+B$.  The rational divisor of the section $s$, defined by
divisible powers, represents $H$.  Consequently
\begin{equation}\label{cone:fano-quotient}
 K_E+B\sim_{\Q}-aH.
\end{equation}
Thus $E$ is of Fano type, the pair has standard coefficients, and its
negative log canonical class is nef.  Birkar's bounded complement
theorem \cite[Theorem~1.7]{BirkarComplements} supplies an lc complement
$B^+\geq B$ with $N(K_E+B^+)\sim0$ for bounded $N$.  For standard
coefficients the complement rounding inequality ensures the stated
monotonicity.

Choose a rational $N$-pluri-volume form on $E$ with divisor $-NB^+$ and
take its logarithmic product $\psi$ with $(d\log s)^N$.  It is lc on
the function field: use the compactification of the torus with both
boundary points.  The same slice ramification calculation makes its
divisor on the punctured cone nonpositive.  Therefore
$f_*=\theta^N/\psi$ has no pole there, is homogeneous of degree $aN$,
and extends across the vertex by normality and $n\geq2$.  Its
discrepancy identity proves \eqref{cone:complement-section}.

Normalize one component of the cover obtained by adjoining
$u^{aN}=f_*$.  Write its ring as $R'$, grade $u$ in degree one, and set
\begin{equation}\label{cone:slice-data}
 \widehat E=\Proj R',\qquad D=aH_u,
\end{equation}
where $H_u$ is the effective rational divisor of $u$, defined using
divisible powers.  A connected grading torus preserves each component,
so the chosen component remains graded.  The variety $\widehat E$ is
normal and projective, and $D$ is ample, effective and $\Q$-Cartier.
Because $u$ has weight one,
$k(\Spec R')=k(\widehat E)(u)$.  Define a rational volume form $\eta$
on $\widehat E$ by the pullback identity
\begin{equation}\label{cone:slice-form}
 \theta=u^a\eta\wedge d\log u.
\end{equation}
The $N$-th power of $\eta\wedge d\log u$ is the pullback of $\psi$,
so $\eta$ is lc.  On $\widehat E\setminus\Supp D$ it has zero Weil
divisor and is klt.  Indeed the cone cover is \'etale where $u\ne0$,
and that open is the product of the corresponding base open with a
torus.

With $h=\vol(v)$ one has
\begin{equation}\label{cone:volume}
 h=a^n H^{n-1},\qquad D^{n-1}\leq Nh.
\end{equation}
For the first equality, count graded pieces up to integral degree $ap$.
Divisible degrees have the ample Hilbert polynomial on $E$.  Every
other degree class has the same leading growth: nonzero degrees generate
$\Z$, so multiplication by fixed homogeneous elements compares each
class above and below with divisible classes, with bounded degree
shifts.  These shifts are fixed for the individual cone and need not be
uniform along the sequence.  Summing yields the first equality.
For the second, the cone cover has degree at most $aN$, and its degree
on Proj is the same because the two homogeneous fraction fields have
weight-one torus parameters.  Since $H_u$ pulls back $H$,
$a^{n-1}H_u^{n-1}\leq a^{n-1}(aN)H^{n-1}=Nh$.
